% Generated by Listas/make_statements.py from pratica-lecture-02.md, item 20. Do not edit.
\textbf{Source.} Tadelis (2013), \emph{Game Theory: An Introduction}, Exercise 5.12 ``Asymmetric Bertrand'', p.~98 {[}PDF p.~117{]}. Answer in the source: none (the preface, p.~xiii, says about half of the exercises are solved on Princeton UP's website). Reproduced verbatim.



\begin{quote}

\textbf{Asymmetric Bertrand:} Consider the Bertrand game with \(c_1(q_1) = q_1\) and \(c_2(q_2) = 2q_2\) and demand equal to \(p = 100 - q\), in which firms must choose prices in increments of \$0.01. We have seen in Section 5.2.4 that one possible Nash equilibrium is \((p_1^*, p_2^*) = (1.99, 2.00)\).



\begin{enumerate}

\def\labelenumi{\alph{enumi}.}

\item

  Show that there are other Nash equilibria for this game.

\item

  How many Nash equilibria does this game have?

\end{enumerate}

\end{quote}



\textbf{Referred to in the statement.} The Bertrand game of Tadelis §5.2.4 (pp.~88--91 {[}PDF p.~107--110{]}) (reproduced from the book):



\begin{quote}

Demand is given by \(p = 100 - q\); each firm posts a price \(p_i\) for otherwise identical goods, all buyers buy from the firm whose price is lowest, and if there is a tie the market splits equally between the two firms. In the asymmetric example \(c_1 = 1\) and \(c_2 = 2\) and prices are limited to increments of \(\epsilon = 0.01\); the book then says that ``one possible Nash equilibrium (you are asked to find more in exercise 5.12 at the end of this chapter) is \((p_1, p_2) = (1.99, 2.00)\).''

\end{quote}

